\section{Where GFM Diverges from Utility Theory}
\label{sec:structural_predictions}

The fungibility-collapse correspondence of
Section~\ref{sec:fungibility_collapse} establishes the regime of
agreement between GFM and standard utility theory.  This section
catalogues the regime of divergence.  Five classes of agent
decisions, generated by combinations of the structural features
of Section~\ref{sec:fc_features}, where GFM makes predictions
that utility theory either does not, or can only match by adding
ad hoc structure that is itself unfalsifiable.  The predictions
are what becomes falsifiable when capability welfare economics is
quantified into a measure-theoretic apparatus: each is testable
in principle through the revealed-sacrifice channel
of~\citep[Revealed~Sacrifice]{lasser2026paper8a} with
interpretation supplied by~\citep[Need~Sufficiency]{lasser2026paper8b}'s
diagnostic architecture.  We do not perform the empirical tests
in this paper; the contribution at this stage is to identify the
predictions clearly enough that they can be tested.

\subsection{Observability prerequisite}
\label{sec:observability_prereq}

Before stating the predictions, a precondition note.  The default
public-ledger-only operation of~\citep[Revealed~Sacrifice]{lasser2026paper8a}
exposes only \emph{event-level} sacrifice magnitudes and does not in
general reveal the agent's full capability inventory or per-capability
cooperative gains.  Several predictions below require richer
observability than the public-ledger default supplies:
\begin{itemize}
\item \textbf{Capability-inventory exposure.}  Predictions 1 and 2
  require knowing whether the agent's pre-trade poset already
  contains \emph{substitutes or complements} relevant to the
  acquired bundle's cooperative components: substitutes that
  reduce marginal value (Prediction~1) or complements that
  amplify it (Prediction~2).
  Operationally this requires either consented-perimeter
  attestation~(\citep[Revealed~Sacrifice]{lasser2026paper8a},
  Appendix~A on the OI--exercise bridge) or hedonic-disaggregation
  panels in which agents' capability inventories are surveyed
  alongside their trades.
\item \textbf{Per-capability sacrifice attribution.}  Prediction~5
  requires decomposing a bundle-for-bundle trade into per-capability
  contributions to the surrendered and acquired sides.  This is the
  bundle-decomposition extension of~\citep[Revealed~Sacrifice]{lasser2026paper8a},
  Remark~5 and Open~Question~2; it is not exposed by the default
  observation protocol.
\item \textbf{Pre-trade need-attestation.}  Prediction~5's two-regime
  bifurcation depends on whether the agent's pre-trade state is
  above or below sufficiency on the relevant dimensions.  This is
  S1-attestation in the sense of~\citep[Need~Sufficiency]{lasser2026paper8b},
  Definition~6, and the discussion of S1-admissibility in
  Section~5 of that paper; it requires institutional infrastructure
  beyond the trade ledger itself.
\end{itemize}
Predictions stated in the public-ledger-only regime require
additional softening.  We note inline where the predictions degrade
under reduced observability, and Section~\ref{sec:empirical}
discusses which deployments naturally support which observability
regimes.

\subsection{Prediction 1: Bundle-completion magnitude}
\label{sec:pred_bundle_completion}

\begin{prediction}[Bundle-completion magnitude]
\label{pred:bundle_completion}
Consider two agents $i$ and $j$ purchasing the same priced bundle
$Y$ for the same nominal price.  Suppose: (a) agent $i$'s
pre-trade capability poset $P_i$ contains \emph{no} substitutes
for $Y$'s cooperative components, while agent $j$'s pre-trade
poset $P_j$ contains substitutes for at least one component of
$Y$ at operationally significant magnitude; and (b) the
deployment satisfies a \emph{bounded-wedge competitive-market
condition}: there exists a deployment-specific bound
$\bar\eta \geq 0$ such that for any two agents $a, b$ with
$\Ui[a](Y) > \Ui[b](Y)$,
\[
  \mathrm{Sacrifice}_a(Y) - \mathrm{Sacrifice}_b(Y)
  \;\geq\;
  \bigl(\Ui[a](Y) - \Ui[b](Y)\bigr) - \bar\eta,
\]
i.e., the cross-agent sacrifice ordering tracks the cross-agent
valuation ordering up to a bounded wedge $\bar\eta$.  Then in
revealed-sacrifice units (the time-or-money cost the agent
surrenders relative to its outside-option benchmark), agent $i$
pays \emph{more} than agent $j$ by at least the
cooperative-gain gap minus the wedge:
\[
  \mathrm{Sacrifice}_i(Y) - \mathrm{Sacrifice}_j(Y)
  \;\geq\;
  \bigl(\Delta\volL_{\mathrm{coop}}(Y \mid P_i)
       - \Delta\volL_{\mathrm{coop}}(Y \mid P_j)\bigr)
  - \bar\eta.
\]
\end{prediction}

\paragraph{Utility-theory alternative.}
A standard-economics treatment matches the prediction by asserting
$\mathrm{MU}_i(Y) > \mathrm{MU}_j(Y)$, i.e., agent~$i$'s marginal utility
for the bundle exceeds agent~$j$'s.  This is a post-hoc utility
assignment, not a prediction: any difference in willingness-to-pay
can be rationalized by the right marginal-utility assignment, and
nothing in standard utility theory constrains $\mathrm{MU}_i$ from
the agent's observable capability inventory.  The
``prediction'' under utility theory is therefore unfalsifiable in
principle: any observation is consistent with some assignment of
$\mathrm{MU}$.

\paragraph{GFM prediction.}
The cooperative-capability term in axiom~M6
(\citep[Proposition~1]{lasser2026poset}) is structurally determined
by the agent's poset.  An observer with consented-perimeter access
to $P_i$ and $P_j$ can compute the cooperative gain
$\Delta\volL_{\mathrm{coop}}(Y \mid P)$ that $Y$ contributes to $P$
on top of its own capability magnitude.  This is the poset-coherent
cooperative excess that M6 admits.  By construction,
$\Delta\volL_{\mathrm{coop}}(Y \mid P_i) >
\Delta\volL_{\mathrm{coop}}(Y \mid P_j)$ when $P_i$ has no
substitutes for $Y$'s components and $P_j$ does.

The bridge from the cooperative-gain ordering to the
sacrifice-magnitude ordering uses the structural reading of
revealed preference~\citep{samuelson1938note} together with the
upper-bound structure of S0.  Under standard
revealed-preference, the revealed-sacrifice bound from
\citep[Theorem~2]{lasser2026paper8a} gives
$\Delta\volL(X) \leq \Ui(Y) \leq \volRwin(Y)$ on the surrendered
side; the agent's valuation $\Ui(Y)$ is bounded below by the
sacrifice and above by the operational truth.  In a competitive
trade environment where agents pay close to their willingness
(an empirical regularity, not a theorem of the framework), an
agent with higher cooperative gain has higher $\Ui(Y)$ and hence
higher willingness to surrender.  The prediction is therefore a
\emph{comparative-statics} claim about cross-agent
willingness-to-pay variance: the sacrifice ordering tracks the
cooperative-gain ordering up to the deployment's competitive-
market gap.

\begin{remark}[Threshold-conditional sign agreement]
\label{rem:pred1_comparative}
The bounded-wedge condition delivers the cooperative-gain $\to$
sacrifice ordering only when the cooperative-gain gap exceeds
the wedge: when
$\Delta\volL_{\mathrm{coop}}(Y \mid P_i)
- \Delta\volL_{\mathrm{coop}}(Y \mid P_j) > \bar\eta$,
Prediction~\ref{pred:bundle_completion}'s lower bound is strictly
positive and the prediction $\mathrm{Sacrifice}_i >
\mathrm{Sacrifice}_j$ holds.  When the cooperative-gain gap is
smaller than $\bar\eta$, the relative-ordering prediction is
empirically indeterminate: differences in observed sacrifices on
that scale are within the deployment's market wedge and do not
support a sign claim either way.  The empirical operability of
the prediction therefore depends on the relative magnitudes of
the cooperative-gain gap and the wedge.  Deployments with small
$\bar\eta$ (competitive markets, low-friction trading) admit
finer-grained tests; deployments with large $\bar\eta$
(monopsony, behavioural anomalies, severe liquidity constraints)
admit only coarse tests.  Standard utility theory cannot
generate even the threshold-conditional comparative-statics
claim structurally, because $\mathrm{MU}_i$ is unconstrained by
the agent's observable capability inventory.
\end{remark}

The prediction is testable on the threshold-restricted regime:
observe $P_i, P_j$ and an estimate or upper bound on $\bar\eta$;
predict $\mathrm{Sacrifice}_i > \mathrm{Sacrifice}_j$ for those
pairs whose cooperative-gain gap exceeds $\bar\eta$.  Violation
of the ordering on this regime falsifies either $\volL$'s
cooperative-excess structure (M6) or the bounded-wedge condition
itself (the deployment's market is less competitive than the
wedge estimate suggested).  Pairs with cooperative-gain gap
within $\bar\eta$ are excluded from the test.

\subsection{Prediction 2: Cooperative premium}
\label{sec:pred_coop_premium}

\begin{prediction}[Cooperative premium]
\label{pred:coop_premium}
Across a population of agents with similar priced-endowment levels,
let $\mathrm{Var}(\mathrm{Sacrifice}(Y) \mid Y \text{ is cooperative})$
denote the cross-sectional variance of revealed-sacrifice magnitudes
for bundles $Y$ with strong cooperative structure (high
$\Delta\volL_{\mathrm{coop}}$ relative to additive content), and
$\mathrm{Var}(\mathrm{Sacrifice}(Y) \mid Y \text{ is non-cooperative})$
the corresponding variance for bundles with little cooperative
content.  Then
\[
  \mathrm{Var}(\mathrm{Sacrifice}(Y) \mid Y \text{ is cooperative})
  \;>\;
  \mathrm{Var}(\mathrm{Sacrifice}(Y) \mid Y \text{ is non-cooperative})
\]
at the population level, after controlling for priced-endowment
heterogeneity.
\end{prediction}

\paragraph{Utility-theory alternative.}
Standard utility theory has no mechanism to differentiate the two
variances structurally.  Both sides of the inequality are
willingness-to-pay variances over a population, and standard utility
theory predicts willingness-to-pay variance from the population's
preference heterogeneity.  Without additional ad hoc structure
asserting that ``cooperative bundles induce more preference
heterogeneity than non-cooperative bundles'' (which is not
derivable from any utility-theoretic axiom), the framework
cannot generate
the directional prediction.  As before, post-hoc utility
assignments can match any observed inequality, but the inequality
is not predicted.

\paragraph{GFM prediction.}
Cooperative value is poset-context-dependent: the cooperative
excess $\Delta\volL_{\mathrm{coop}}(Y \mid P)$ varies with the
agent's poset $P$ in ways that the additive content
$\sum_{c \in Y} w(c)$ does not (the additive content is the same
for any agent with the same priced-endowment level).  Across a
population with heterogeneous posets, the cooperative bundles
inherit poset-context heterogeneity that the non-cooperative
bundles do not, and the revealed-sacrifice variance reflects the
inherited heterogeneity.  The prediction is testable through
cross-sectional sacrifice-magnitude variance comparisons after
controlling for priced-endowment and demographic confounds.

\begin{remark}[Operational caveat]
\label{rem:coop_premium_op}
Prediction~\ref{pred:coop_premium} is robust to the
public-ledger-only observability regime in the sense that
revealed-sacrifice variance is computable from event-level data
without per-agent capability-inventory exposure, provided the
population can be stratified into ``cooperative-heavy'' versus
``non-cooperative-heavy'' bundle categories ex ante.  The
stratification itself, however, requires structural knowledge of
which bundles support cooperative excess, a content question that
\citep[Need~Sufficiency]{lasser2026paper8b}'s admissibility
discipline (Definition~7 alarm condition, and the boundary
construction in Section~5) is positioned to inform.
\end{remark}

\subsection{Prediction 3: Saturation transition}
\label{sec:pred_saturation}

\begin{prediction}[Saturation transition]
\label{pred:saturation}
For an agent whose capability poset $P$ becomes saturated in a
dimension $k$ (further additions to $k$ contribute negligibly to
$\volL(P)$ at the operationally relevant scale), the agent's
revealed-sacrifice pattern over the saturating window exhibits a
\emph{transition}: revealed sacrifice toward dimension $k$ declines
toward a structural floor as saturation approaches, while
revealed sacrifice toward unsaturated dimensions does not decline
in the same direction.  Longitudinally, the transition occurs at
the saturation point in $k$ that the poset's structural
connectivity predicts, not at an exogenously assigned utility-curve
inflection.
\end{prediction}

\paragraph{Utility-theory alternative.}
Standard utility theory under concave $u$ predicts diminishing
sacrifice magnitudes as wealth in any dimension grows, but does
not predict \emph{which} dimensions saturate or \emph{when} the
transition occurs from a structural ground.  The agent's curvature
of $u$ in dimension $k$ is a free parameter; given any saturation
trajectory in observed sacrifice, a $u$-shape can be fit to match.
The prediction is descriptive rather than structural.

\paragraph{GFM prediction.}
The saturation point in dimension $k$ is determined by the poset's
structural connectivity in that dimension: specifically, by the
$\delta_{\mathrm{covered}}$ contribution to the cell decomposition
of \citep[Need~Sufficiency]{lasser2026paper8b}, Definition~8, on
the dimension's downstream cone.  When the downstream cone is
exhausted (every realized output that can be produced from
dimension $k$ \emph{is} being produced at the agent's current
$\volR$), additional capability in $k$ contributes only to
$\delta_{\mathrm{dormant}}$.  This is observable from the agent's
production trajectory: the saturation point coincides with the
exhaustion of dimension $k$'s downstream cone in the agent's
realized output set.  Standard utility theory has no
counterpart: the curvature of $u$ in $k$ is a parameter, not a
structural prediction from the agent's production.

\begin{remark}[Saturation as a longitudinal signature]
\label{rem:saturation_longitudinal}
The prediction is fundamentally longitudinal.  Cross-sectionally,
two agents at the same wealth level may differ in their
saturation status in dimension $k$ depending on their poset
compositions; the sacrifice pattern at any single time is
therefore an unreliable saturation diagnostic.  The signature
that distinguishes the GFM and utility-theoretic predictions is the
\emph{trajectory}: GFM predicts the transition occurs at the
saturation point implied by the downstream-cone exhaustion;
utility theory predicts it occurs at the curvature inflection of
$u$, which is not anchored to observable production structure.
\end{remark}

\subsection{Prediction 4: Anti-monopolar at population scale}
\label{sec:pred_anti_monopolar}

\begin{prediction}[Anti-monopolar at population scale]
\label{pred:anti_monopolar}
Across populations $\Pi_1, \dots, \Pi_M$ with comparable
aggregate capability magnitudes $\volL(\Pi_m) \approx \bar V$ but
heterogeneous diversity (measured via revealed-sacrifice dispersion
across capability categories), populations with higher diversity
exhibit higher long-run welfare correlates than populations with
equivalent total capability magnitude but lower diversity.
Operationally, the population-scale wireheading-consistent HHI
of~\citep[Need~Sufficiency]{lasser2026paper8b}, Proposition~6,
computed cross-sectionally on trade-flow concentration, correlates
\emph{negatively} with longitudinal welfare-correlate trajectories.
\end{prediction}

\paragraph{Utility-theory alternative.}
Standard utility theory aggregated to the population level (via,
e.g., representative-agent reasoning or income-weighted welfare
indices) is silent on diversity.  Welfare under linear
representative-agent utility depends on aggregate wealth; under
concave representative-agent utility, on aggregate wealth and its
distribution; but in either case, the \emph{categorical
distribution} of capabilities (rather than the wealth
distribution) does not enter the welfare calculation.  Predictions
about cross-population welfare from diversity require leaving
standard utility theory and adopting a multi-attribute welfare
function whose attribute-mixture structure is itself unconstrained.

\paragraph{GFM prediction.}
\citep[Proposition~6]{lasser2026horizon} establishes the
anti-monopolar property at the agent / sub-collective level: full
domination by a single substrate is anti-maximizing under
discounting.  The population-scale claim of
Prediction~\ref{pred:anti_monopolar} is the cross-population
analogue: populations whose capability exercise concentrates in a
narrow set of dimensions are exposed to the discounting-induced
welfare loss the anti-monopolar property identifies, and the
wireheading-consistent HHI provides the operational signal for
detecting concentration.  The prediction is testable cross-nationally
through long-run welfare-correlate panels (life satisfaction,
inequality-adjusted human-development indices, multi-dimensional
poverty measures) regressed against trade-flow HHI from
revealed-sacrifice data.

\begin{remark}[Concentration consistent with, not causal proof]
\label{rem:hhi_consistent_with}
The wireheading-consistent HHI is a \emph{deployment-readable
signal} for concentration in trade-flow data; it is not by itself
a causal-inference instrument.  Concentration consistent with
anti-monopolar welfare loss may also reflect market consolidation,
demand shocks, or supply-side constraints that have nothing to do
with the anti-monopolar mechanism.  Prediction~\ref{pred:anti_monopolar}
asserts a correlational signature consistent with the
anti-monopolar mechanism; isolating the mechanism from
alternatives requires the kind of natural-experiment or instrumental-
variable design that empirical economics deploys for related
welfare questions.  The framework's contribution is the
deployment-readable signal, not a causal-identification strategy.
\end{remark}

\subsection{Prediction 5: Mixed-polarity bundle-for-bundle traces}
\label{sec:pred_mixed_polarity}

\begin{prediction}[Mixed-polarity bundle-for-bundle traces]
\label{pred:mixed_polarity}
Consider a bundle-for-bundle trade (an agent surrenders a bundle
$X = \{x_1, \dots, x_p\}$ for an acquired bundle
$Y = \{y_1, \dots, y_q\}$ at the same nominal price) in which the
agent's pre-trade state is at or below sufficiency
(\citep[Need~Sufficiency]{lasser2026paper8b}, Definition~6) on at
least one dimension covered by $Y$.  Two agents performing
nominally identical priced trades produce different need-share
versus want-share splits in the bundle decomposition, depending
on their pre-trade need-satisfaction state:
\[
  \frac{\mathrm{NAC}_i(Y)}{\mathrm{Sacrifice}_i(X)}
  \;\neq\;
  \frac{\mathrm{NAC}_j(Y)}{\mathrm{Sacrifice}_j(X)},
\]
where $\mathrm{NAC}_i(Y)$ is the need-access-cost contribution
of~\citep[Need~Sufficiency]{lasser2026paper8b}, Definition~9,
attributable to the need-satisfying components of $Y$ for agent~$i$.
The split bifurcates the population into two regimes (high-NAC,
low-NAC) on otherwise identical priced trades.
\end{prediction}

\paragraph{Utility-theory alternative.}
Standard utility theory observes the priced bundle-for-bundle trade
as a single utility delta $\Delta U = U(Y) - U(X) > 0$ (the agent
chose $Y$).  No further structural decomposition is available;
the polarity of the trade is collapsed into a scalar utility
change.  Two agents performing the same priced trade are, under
linear utility, treated identically; under heterogeneous concave
$u_i$, the variance is attributed to $u_i$'s shape rather than to
the agents' pre-trade need-satisfaction states.  The bifurcation
prediction is therefore not generated by standard utility theory.

\paragraph{GFM prediction.}
\citep[Need~Sufficiency]{lasser2026paper8b} establishes the
polarity boundary (Definition~3) and the need-access cost
(Definition~9) as separate measurement objects from the want-side
revealed sacrifice.  Below the polarity boundary, sacrifice
reveals cost-of-access (the agent had no non-degrading
alternative); above, sacrifice reveals preference (the agent could
have done otherwise).  When a single trade carries components
on both sides of the polarity boundary, as bundle-for-bundle
trades canonically do (residential relocation, discussed at length
in Section~\ref{sec:worked_examples}), the mixed polarity
bifurcates the population into NAC-dominant versus
above-sufficiency-dominant regimes on identical priced trades.
The prediction requires the bundle-decomposition extension of
\citep[Revealed~Sacrifice]{lasser2026paper8a},
Remark~5 and Open~Question~2, and is therefore conditional on
that extension being formalized.

\begin{remark}[The prediction is the cleanest test]
\label{rem:pred5_clean_test}
Prediction~\ref{pred:mixed_polarity} is the structurally cleanest
test of the GFM-vs-utility-theory divergence in the catalogue.
The other four predictions can be matched by utility theory at the
cost of post-hoc structure (heterogeneous $u_i$, multi-attribute
welfare functions, non-stationary curvature in $u$); the bifurcation
on identical priced trades cannot, because standard utility theory
treats identical priced trades identically by construction.  The
deployment cost is correspondingly higher: it requires bundle
decomposition, pre-trade need-attestation, and S1-attestation
infrastructure that the public-ledger default does not supply.  But
when the deployment supports the test, the prediction is sharp.
\end{remark}

\subsection{Summary table}
\label{sec:pred_summary}

\begin{table}[h]
\centering
\small
\begin{tabular}{p{0.20\linewidth} p{0.25\linewidth} p{0.20\linewidth} p{0.20\linewidth}}
\toprule
\textbf{Prediction} & \textbf{Mechanism} &
\textbf{Observability} & \textbf{Falsification} \\
\midrule
P1: Bundle-completion magnitude &
Cooperative excess scales with poset gap &
Capability-inventory + sacrifice ledger &
Falsifies $\volL$ \emph{or} S0 if violated \\
\addlinespace
P2: Cooperative premium &
Poset-context heterogeneity inflates variance &
Stratifiable bundle ledger &
Falsifies M6's directional content \\
\addlinespace
P3: Saturation transition &
Downstream-cone exhaustion sets transition point &
Longitudinal sacrifice + production traces &
Falsifies the cone-exhaustion mechanism \\
\addlinespace
P4: Anti-monopolar at population scale &
Cross-population diversity correlates with welfare &
Population HHI + welfare panels &
Falsifies the population-scale extension \\
\addlinespace
P5: Mixed-polarity bundle-for-bundle &
Pre-trade need state bifurcates NAC share &
Bundle decomp + S1-attestation &
Falsifies polarity-boundary discipline \\
\bottomrule
\end{tabular}
\caption{Summary of the five structural predictions.  Each row
identifies the GFM mechanism, the observability regime required to
test the prediction, and what a violation of the prediction would
falsify.  None of the falsification routes overlap, so independent
empirical tests of the predictions provide independent evidence
for the structural-mechanism claims.}
\label{tab:predictions}
\end{table}

The five predictions span all five structural features of
Section~\ref{sec:fc_features}: P1 and P2 lean on cooperative
capabilities (M6); P3 leans on the downstream cone and the
exercise-pathway machinery of \citep[Horizon~Aware]{lasser2026horizon}
re-cast through \citep[Need~Sufficiency]{lasser2026paper8b}'s
$\delta$-decomposition; P4 leans on the anti-monopolar property; P5
leans on individuation, polarity, and the bundle-decomposition
extension simultaneously.  No single feature alone generates all
five predictions; the structural-prediction class is a property of
the apparatus as a whole, not of any one axiom.  The
fungibility-collapse correspondence of
Section~\ref{sec:fungibility_collapse} explains why utility theory
cannot generate the predictions structurally: under fungibility
collapse, the mechanism that each prediction relies on degenerates
to the single-dimension wealth axis, and the directional content of
the prediction vanishes.
